\documentclass[10pt,a4paper]{amsart}
\usepackage[margin=19mm]{geometry}
\usepackage{amsmath,amssymb,mathtools}
\usepackage[scr=rsfs,cal=euler]{mathalfa}
\usepackage{xfrac}
\usepackage[unicode,hidelinks,bookmarks=false]{hyperref}
\newcommand{\ie}{i.e.\ }
\newcommand{\cf}{cf.\ }
\newcommand{\resp}{respectively\ }
\newcommand{\Subsection}{\subsection}
\providecommand{\nobreakdash}{\nobreak\hbox{-}}
\setlength{\emergencystretch}{2em}
\multlinegap0pt
\usepackage{amssymb}
\usepackage{float}
\usepackage{bm}
\usepackage{cancel}
\usepackage{xcolor}
\newtheorem{theorem}{Theorem}
\newcommand{\calA}{\mathcal{A}}
\title{Generalizations of the Squircle-Lemniscate Relation and Keplerian Dynamics}
\author{Zbigniew Fiedorowicz, Muthu Veerappan Ramalingam}
\date{Source: arXiv:2601.17358v5 (CC BY 4.0)}
\begin{document}
\maketitle

\noindent\emph{Source and licence.} Generalizations of the Squircle-Lemniscate Relation and Keplerian Dynamics. Version arXiv:2601.17358v5 (CC BY 4.0), distributed by arXiv under the Creative Commons Attribution 4.0 International licence, http://creativecommons.org/licenses/by/4.0/, as recorded in the arXiv licence metadata for this exact version and shown on the abstract page https://arxiv.org/abs/2601.17358v5. Full licence notice: \texttt{licenses/arxiv-2601.17358-CC-BY-4.0.txt}.\par\smallskip

\noindent\textit{Editorial scope.} Counted unit: the article's theorem fundamental\_integral\_theorem (source0.tex lines 163--168). The article's complete original proof of it is delivered with it (source0.tex lines 171--213, one \texttt{proof} environment of 43 lines). No mathematics is written, repaired or re-derived: the statement and the proof are the article's own lines, in the article's own notation, and no retained line is substituted or re-flowed.  No further source-local context is needed: every cross-reference of every retained line resolves inside the two delivered spans.  Nothing was omitted from any retained span, so the package carries no omission record.  No retained line contains a citation, so the wrapper carries no bibliography and no bibliography entry.  No retained line contains a figure environment, an image-inclusion command or drawing code, so no image is delivered and the package declares no asset dependency.  Editorial formatting only, in addition: an AMS \texttt{amsart} preamble that restates the article's own declarations and macros used by the retained lines, namely the environments \texttt{theorem} on separate counters, together with the standard packages those lines need.  The retained environments print in this document as Theorem 1 in order of appearance, whereas the article prints the same items with the numbering of its own full document.  The complete native source of the article as distributed in the arXiv e-print for this exact version is bundled unchanged as \texttt{sources/193310/source0.tex}.
\begin{theorem}\label{fundamental_integral_theorem}
For any positive integer $n$, the following equality holds:
\begin{equation}\label{fundamental_identity}
    \int_0^1 \frac{\,dr}{\sqrt{1-r^{2n}}} = 2^{1/n} \int_0^1 \sqrt[2n]{1-x^{2n}} \, dx
\end{equation}
\end{theorem}

\begin{proof}
The integral $\calA=\int_0^1 \sqrt[2n]{1-x^{2n}} \, dx$ represents the area of the portion of
the Lam\'e curve $x^{2n}+y^{2n}=1$ contained within the first quadrant. An alternative way of
computing this area is to use the following parametrization of the Lam\'e curve:
\[x=\cos^{\frac{1}{n}}(nt),\ y=\sin^{\frac{1}{n}}(nt),\quad 0\le t\le\frac{\pi}{2n}.\]
(Note that these are the same as the polar equations of the sinusoidal spiral restricted to the interval $\left[0,\frac{\pi}{2n}\right]$.)
Thus by Green's theorem, we have
\begin{align*}
\calA &=\frac{1}{2}\int_{\partial\{(x,y)\,|\,x^{2n}+y^{2n}\le1,\,x\ge0,\,y\ge0\}}(x\,dy-y\,dx)\\
&=0+0+\frac{1}{2}\int_0^{\frac{\pi}{2n}}\left(x\frac{dy}{dt}-y\frac{dx}{dt}\right)\,dt\\
&=\frac{1}{2}\int_0^{\frac{\pi}{2n}}\left(\frac{\cos^{\frac{1}{n}}(nt)\sin^{\frac{1}{n}}(nt)\cos(nt)}{\sin(nt)}+
\frac{\sin^{\frac{1}{n}}(nt)\cos^{\frac{1}{n}}(nt)\sin(nt)}{\cos(nt)}\right)\,dt\\
&=\frac{1}{2}\int_0^{\frac{\pi}{2n}}\left(\frac{\cos^{\frac{1}{n}}(nt)\sin^{\frac{1}{n}}(nt)\left(\cos^2(nt)+\sin^2(nt)\right)}{\cos(nt)\sin(nt)}\right)\,dt\\
&=\frac{1}{2^{\frac{1}{n}}}\int_0^{\frac{\pi}{2n}}\left(\frac{2^{\frac{1}{n}}\cos^{\frac{1}{n}}(nt)\sin^{\frac{1}{n}}(nt)}{2\cos(nt)\sin(nt)}\right)\,dt\\
&=\frac{1}{2^{\frac{1}{n}}}\int_0^{\frac{\pi}{2n}}\frac{\sin^{\frac{1}{n}}(2nt)}{\sin(2nt)}\,dt
\end{align*}
Substituting $u=2t$ into the last integral, we obtain
\begin{align}
\calA&=\frac{1}{2^{1+\frac{1}{n}}}\int_0^{\frac{\pi}{n}}\frac{\sin^{\frac{1}{n}}(nu)}{\sin(nu)}\,du\nonumber\\
&=
\frac{1}{2^{1+\frac{1}{n}}}\left[\int_0^{\frac{\pi}{2n}}\frac{\sin^{\frac{1}{n}}(nu)}{\sin(nu)}\,du+\int_{\frac{\pi}{2n}}^{\frac{\pi}{n}}\frac{\sin^{\frac{1}{n}}(nu)}{\sin(nu)}\,du\right]\label{area_eqn2}
\end{align}
Using the substitution $v=\frac{\pi}{n}-u$, we see that
\begin{align*}
\int_{\frac{\pi}{2n}}^{\frac{\pi}{n}}\frac{\sin^{\frac{1}{n}}(nu)}{\sin(nu)}\,du &=-\int_{\frac{\pi}{2n}}^0\frac{\sin^{\frac{1}{n}}(\pi-nv)}{\sin(\pi-nv)}\,dv\\
&=\int_0^{\frac{\pi}{2n}}\frac{\sin^{\frac{1}{n}}(nv)}{\sin(nv)}\,dv\\
&=\int_0^{\frac{\pi}{2n}}\frac{\sin^{\frac{1}{n}}(nu)}{\sin(nu)}\,du
\end{align*}
Thus \eqref{area_eqn2} simplifies to
\begin{equation}
\calA=\frac{1}{2^{1+\frac{1}{n}}}\left[2\int_0^{\frac{\pi}{2n}}\frac{\sin^{\frac{1}{n}}(nu)}{\sin(nu)}\,du\right]=
\frac{1}{2^{\frac{1}{n}}}\int_0^{\frac{\pi}{2n}}\frac{\sin^{\frac{1}{n}}(nu)}{\sin(nu)}\,du.
\end{equation}
Making the substitution $r=\sin^{\frac{1}{n}}(nu)$, we obtain
\[\frac{dr}{du}=\frac{\sin^{\frac{1}{n}}(nu)\cos(nu)}{\sin(nu)}=\frac{\sin^{\frac{1}{n}}(nu)\sqrt{1-r^{2n}}}{\sin(nu)},\]
which implies that
\begin{equation}
\calA=\frac{1}{2^{\frac{1}{n}}}\int_0^1\frac{\,dr}{\sqrt{1-r^{2n}}}
\end{equation}
Hence
\[\int_0^1\frac{\,dr}{\sqrt{1-r^{2n}}}=2^{\frac{1}{n}}\calA=2^{\frac{1}{n}}\int_0^1 \sqrt[2n]{1-x^{2n}} \, dx.\]
This completes the proof.
\end{proof}

\end{document}
